\section{Operator registry}
\label{app:operators}

Table~\ref{tab:operators} is the full registry of the \ConfigValue{NUM_ENABLED_FAMILIES} production
operator families (of \ConfigValue{NUM_MARK_FAMILIES_TOTAL} defined; \ConfigValue{DISABLED_FAMILIES} are
disabled everywhere and enumerate no candidates). Every family's rewrite is lowered to validity-checked
primitives; the ``successor invariant'' column states what the family guarantees beyond the universal
connected-and-valence-valid closure of \Cref{prop:closure}. Charge is preserved by every family (there is
no charge-changing transition). ``Training source'' names which training layer typically supervises the
family (\Cref{app:data}); ring-topology growth is retained from the de novo carbon-tree prior because the
edit layers preserve ring count.

\begin{table}[h]
\centering
\caption{Production operator families. All families preserve connectivity, per-atom valence, and net
charge (\Cref{prop:closure}); implicit hydrogens re-derive from each touched atom's valence class.}
\label{tab:operators}
\footnotesize
\renewcommand{\arraystretch}{1.15}
\begin{tabularx}{\textwidth}{@{}l Y Y l@{}}
\toprule
Family & Operands / applicability & Successor invariant \& inverse & Training source \\
\midrule
atom\_insert & peripheral site + element/valence class + incident bond order & adds one heavy atom
($\le N$); inverse = atom\_delete & corruption; MMP graft-in \\
atom\_delete & peripheral (leaf) atom whose removal keeps connectivity & removes one heavy atom; inverse =
atom\_insert & corruption; MMP peel \\
atom\_restate & non-ring atom; new valence class feasible at its heavy-bond sum & element fixed, valence
class changed, H re-derived; self-inverse family & corruption (bioisostere) \\
bond\_reorder & non-ring (acyclic) bond; new integer order & bond order changed, endpoints' H re-derived;
self-inverse family & corruption \\
bond\_reroute (Graft) & pendant subtree across a single bridging bond; target attachment site & subtree
relocated, ring core fixed; scored at canonical successor group; inverse = reroute back & corruption;
shared-scaffold \\
ring\_system\_grow & saturated chain match; typed ring template (mono/fused/bridged/spiro; aromatic/sat.)
& adds a whole ring system with per-atom identity; inverse = ring\_system\_delete & retained from de novo
prior \\
ring\_system\_delete & a supported ring system; decoration-preserving open & opens a ring keeping each
atom's element (Kekul\'e); inverse = ring\_system\_grow on the matching scaffold & corruption (ring-open,
trim direction) \\
ring\_system\_restate & a ring system; aromaticity-changing restate (both directions) & changes perceived
aromaticity via a validity-checked bond-order relayout; self-inverse family & corruption
(aromatize$\leftrightarrow$de-aromatize) \\
\bottomrule
\end{tabularx}
\end{table}

\emph{Disabled families.} \ConfigValue{DISABLED_FAMILIES} are legacy null-prior operators retained in the
mark enumeration for historical compatibility but enumerate no candidates and receive no training signal;
ring additions go through \texttt{ring\_system\_grow}. \emph{Charge.} No family alters an atom's formal
charge; charged centers are protected in the corruption and the net charge is invariant under every edit,
so charge is preserved but not designed (\Cref{sec:limitations}). \emph{Aromaticity.} States are Kekul\'e;
``aromatic'' is a perception on the round-tripped molecule, so \texttt{ring\_system\_restate} realizes
aromatization and de-aromatization as validity-checked integer bond-order relayouts rather than edits of a
stored aromatic bond class.
